% =====================================================================
% Mandatory CAISc 2026 checklists. Rendered with the official style-file
% macros (\involvementA, \iterationLow, \answerYes, ...). Instruction
% blocks have been deleted per template; section/subsection headings and
% all questions are kept verbatim.
% =====================================================================

\newpage

\section*{Conference For AI Scientists 2026 - AI Involvement Checklist}
\label{checklist_ai_involvement}

\subsection*{Research Stage Assessment}

\begin{enumerate}

  \item \textbf{Hypothesis development}.

  Answer: \involvementA{}

  Iteration Effort: \iterationNA{}

  Explanation: The research problem, seed idea, and the three-probe
  triangulation framework were conceived and developed by the human authors.
  AI tools were not used to formulate the core research question or the
  framework design.

  \item \textbf{Experimental design and implementation}.

  Answer: \involvementB{}

  Iteration Effort: \iterationHigh{}

  Explanation: The experimental design, choice of baselines, problem families,
  and overall methodology were decided by the human authors. AI coding
  assistants implemented the experimental stack under human direction (the six
  variant-generation scripts, the five-model API sweep harnesses, the three
  family-specific verifiers, and the Infini-gram proximity pipeline) across
  many iterations of debugging and refinement. All design decisions, audit
  sign-off, and final code review were performed by the human authors.

  \item \textbf{Analysis of data and interpretation of results}.

  Answer: \involvementB{}

  Iteration Effort: \iterationMedium{}

  Explanation: Data processing and statistical analysis were primarily carried
  out by the human authors. AI tools assisted with data organisation and
  exploratory checks (for example, reconciling aggregate counts against the
  per-instance verifier logs and surfacing an off-bank duplicate-ID issue that
  the authors then excluded). All reported conclusions were drawn and verified
  by the human authors.

  \item \textbf{Writing}.

  Answer: \involvementB{}

  Iteration Effort: \iterationMedium{}

  Explanation: The manuscript was primarily written by the human authors. AI
  tools were used as writing aids for grammar, sentence restructuring, and
  clarity of exposition. All scientific content, arguments, and conclusions are
  the work of the human authors.

\end{enumerate}

\subsection*{AI System Documentation}

\begin{enumerate}
  \setcounter{enumi}{4}

  \item \textbf{AI systems used}.

  Description: A frontier large language model (accessed through a commercial
  chat interface and API) was used as a coding assistant, writing aid, and
  literature-review aid throughout the project, together with an AI-assisted
  integrated development environment for code completion and refactoring. A
  separate open-weight transformer language model (approximately 7B parameters)
  was the \emph{subject} of the mechanistic analysis in
  Appendix~\ref{app:mech}; it was not used as a research assistant. No custom
  multi-agent frameworks were used. Specific product and model names are
  withheld in this anonymized submission and will be disclosed in the
  camera-ready version.

  \item \textbf{Observed AI Limitations}.

  Description: Two limitations were observed. First, context inconsistency: AI
  coding assistance occasionally produced code subtly inconsistent with the
  existing pipeline (mismatched column names, off-by-one step indices, and in
  one case computing a GSM aggregate over off-bank duplicate problem IDs), each
  caught by human review before integration. Second, imprecise statistical
  language: AI-assisted writing drafts sometimes conflated correlation strength
  with statistical significance, which the human authors corrected before
  inclusion.

\end{enumerate}

\newpage

\section*{Conference For AI Scientists 2026 - Reproducibility and Responsibility Checklist}
\label{checklist_reproducibility}

\begin{enumerate}

\item {\bf Claims}
    \item[] Question: Do the main claims made in the abstract and introduction
    accurately reflect the paper's contributions and scope?
    \item[] Answer: \answerYes{}
    \item[] Justification: The abstract and introduction state three concrete
    findings (block-rename sign flip, WIS universal collapse, and
    injection-compliance/correctness dissociation), each supported in
    Sections~\ref{sec:inversion}--\ref{sec:bw} and~\ref{sec:conclusion} with
    explicit $n$, confidence intervals, and $p$-values. Scope is explicitly
    restricted to zero-shot structured reasoning.

\item {\bf Limitations}
    \item[] Question: Does the paper discuss the limitations of the work
    performed by the authors?
    \item[] Answer: \answerYes{}
    \item[] Justification: A dedicated Limitations section covers scope
    (English, zero-shot, structured reasoning), the Probe-3 proximity-proxy
    status, the Blocksworld measurement-protocol abort, single
    reasoning-trained provider, the behavioural-only scope of CCI, and deferred
    API-budget experiments. An explicit per-(probe, model) coverage matrix is
    given in Appendix~\ref{app:cov}.

\item {\bf Theory assumptions and proofs}
    \item[] Question: For each theoretical result, does the paper provide the
    full set of assumptions and a complete (and correct) proof?
    \item[] Answer: \answerNA{}
    \item[] Justification: The paper is empirical; it presents no formal
    theorems or proofs. All statistical methods and assumptions are documented
    in Appendix~\ref{app:metrics}.

\item {\bf Experimental result reproducibility}
    \item[] Question: Does the paper fully disclose all the information needed
    to reproduce the main experimental results?
    \item[] Answer: \answerYes{}
    \item[] Justification: All prompts and classifier rules are in
    Appendix~\ref{app:prompts}; all metric definitions in
    Appendix~\ref{app:metrics}; all model, temperature, and API settings in
    Section~\ref{sec:setup}. Per-instance verifier logs and a re-derivation
    script (which applies the bank-ID filter so the GSM GPT-4o/Llama figures
    reproduce at $n{=}20$) are shipped in the anonymous repository
    (Appendix~\ref{app:repro}).

\item {\bf Open access to data and code}
    \item[] Question: Does the paper provide open access to the data and code,
    with sufficient instructions to faithfully reproduce the main results?
    \item[] Answer: \answerYes{}
    \item[] Justification: The anonymous repository (Appendix~\ref{app:repro})
    contains per-instance verifier logs, variant-generation scripts, sweep
    drivers, and a re-derivation script, with no author-identifying
    information.

\item {\bf Experimental setting/details}
    \item[] Question: Does the paper specify all the training and test details
    necessary to understand the results?
    \item[] Answer: \answerYes{}
    \item[] Justification: Section~\ref{sec:setup} specifies the five models,
    temperature ($T{=}0$), API provider (OpenRouter), zero-shot CoT prompting,
    and problem-bank sizes; Appendix~\ref{app:pipeline} gives the four-stage
    data flow. No model training is performed.

\item {\bf Experiment statistical significance}
    \item[] Question: Does the paper report error bars suitably and correctly
    defined or other appropriate information about statistical significance?
    \item[] Answer: \answerYes{}
    \item[] Justification: Wilson 95\% confidence intervals for GSM and BW
    proportions; ALGO intervals are 10k percentile cluster bootstrap over
    clone families (seed 42). Paired Wilcoxon with exact $W$ and $p$ for CCI comparisons;
    Pearson $r$ with bootstrap $p$ for correlations; Fisher's exact for
    matched-subset tests. A Holm--Bonferroni multiple-comparison policy is
    documented in Appendix~\ref{app:metrics}.

\item {\bf Experiments compute resources}
    \item[] Question: For each experiment, does the paper provide sufficient
    information on the compute resources needed to reproduce the experiments?
    \item[] Answer: \answerYes{}
    \item[] Justification: The behavioural study comprises approximately 20K
    model API calls via OpenRouter (Table~\ref{tab:coverage}), requiring no
    local GPU. The mechanistic Qwen-2.5-7B pilot was run on a single GPU
    (Google Colab T4), noted in Appendix~\ref{app:mech}.

\item {\bf Research Ethics}
    \item[] Question: Does the research conducted in the paper conform with the
    NeurIPS Code of Ethics, except where CAISc's AI authorship policies apply?
    \item[] Answer: \answerYes{}
    \item[] Justification: The study evaluates publicly available models on
    reasoning benchmarks; there are no human subjects, sensitive data, or
    dual-use risks. AI involvement is documented in the AI Involvement
    Checklist above.

\item {\bf Broader impacts}
    \item[] Question: Does the paper discuss both potential positive and
    negative societal impacts of the work performed?
    \item[] Answer: \answerYes{}
    \item[] Justification: The primary impact is positive---surfacing hidden
    capability gaps that benchmark accuracy obscures supports more informed
    deployment decisions. A potential negative use is adversarially crafting
    benchmark problems that exploit surface fragility; this is mitigated by
    releasing all methods and problem banks so the community can build
    countermeasures.

\end{enumerate}
